\begin{lemma}[The common color choice in Situation 7.3]
Fix a finite loopless multigraph with its individual edge identities, a
palette of \(k\) colors, a distinguished edge \(e_0\), an input coloring,
an exact terminal tree-sequence \(T\), a generating coloring \(\varphi_n\),
a positive-rung extended Tashkinov tree on these data, a hierarchy base, a
hierarchy counterexample \(C\), and a minimum-counterexample context for
\(C\).  Let \(P\) be a
\noderef{HierarchyReplacementPresentation} and write
\(\mu=P.\mathit{coloring}\), \(H=P.\mathit{counterexample.hierarchy}\),
\(\mathcal G=P.\mathit{counterexample.good}\),
\(q=H.\mathit{dividerCount}\), and
\[
 j=\noderef{HierarchyPathNumber}(C.\mathit{hierarchy}),\qquad
 A=\noderef{HierarchyStarredStage}(H,q),\qquad
 \ell=\noderef{HierarchyFullTailLength}(C),\qquad
 B=\noderef{HierarchyFullTailPrefix}(C,j-2).
\]
Let \(T_n\) be the last ladder tree of \(P\)'s replacement ETT and put
\[
 D=D_{n,q}=\noderef{HierarchyStoredColorsAt}(H,q),\qquad
 \Gamma=\bigcup_{\eta\in D}\mathcal G.\Gamma^q_\eta,
\]
\[
 M_0=\bigcup_{v\in
   V(P.\mathit{replacementBase.core})\setminus V(T_n)}
        \overline\mu(v).
\]
For every natural number \(s\), define the finite set of colors on the
displayed tail through index \(s\) by
\[
 U_s=\{c:\text{ for some }r<|T|,\ |A|\le r<|A|+s
                 \text{ and }\mu(T.\mathit{edge}(r))=c\}.
\]
Thus the quantified \(r\) names an exact edge identity of \(T\), not merely
an endpoint pair.

Assume \(2\le j<\ell\).  Let \(\alpha\) be a color which is missing at the
last displayed vertex of \(T\) and belongs to
\(\overline\mu(B)\), and which satisfies the Claim \(j-1\) alternative
\[
 \alpha\notin\Gamma\cup M_0,
 \quad\hbox{or}\quad
 \alpha\in\Gamma^q_\eta\text{ for some }\eta\in D
 \text{ whose first \(\mu\)-missing position is at most }|B|.
\]
Then there is a color \(\beta\) such that
\[
 \beta\ne\alpha,\qquad
 \beta\in\overline\mu(B),\qquad
 \beta\notin\Gamma,\qquad
 \beta\notin U_j,
\]
and one of the following two alternatives holds:
\[
 \beta\notin M_0\cup D,
 \quad\hbox{or}\quad
 \beta\in D\ \text{ and }
   (\mathcal G.\Gamma^q_\beta\cup\{\beta\})\cap U_j=\varnothing.
\]
\end{lemma}
\begin{proof}
Let \(H_0=C.\mathit{hierarchy}\), and let \(A_0\) be its last starred
stage.  The \(\mathit{sameGood}\) field of the replacement presentation is a
\noderef{SameGoodHierarchy}.  Its \noderef{SameHierarchyStages} component
identifies the divider counts, every ordered hierarchy stage, and the
starred bases of \(H_0\) and \(H\).  In particular,
\[
                         A_0=A.                              \tag{1}
\]
These are equalities of tree-sequences, so by
\noderef{TreeSequenceInitialSegment} they preserve the complete order and
every individual displayed edge identity.  From
\noderef{HierarchyFullTailLength}, (1) gives
\[
 |T|-|A|=|T|-|A_0|=\ell.                                    \tag{2}
\]
Likewise, the definition \noderef{HierarchyFullTailPrefix}, together with
\noderef{TreeSequencePrefix}, identifies \(B\) with the replacement
hierarchy's ordered prefix of length
\(\min\{|A|+(j-2),|T|\}\).                                  \tag{3}

We next verify the hypotheses needed for Lemma 9n under the replacement
coloring.  The field \(P.\mathit{minimumContext}\) is precisely a
\noderef{HierarchyMinimumCounterexampleContext} for \(P\)'s replacement
input, ETT, base, counterexample, and coloring \(\mu\).  Supply its
properness, exact-one-uncoloredness, lower-rung package, fewer-divider
induction, lexicographic-minimality, and large-missing-set fields to
\noderef{HierarchyCounterexamplePenultimateElementary}.  The result is
that the replacement prefix at full-tail index \(\ell-1\) is elementary
under \(\mu\).  Here its full-tail length is \(\ell\) by (2), and its
prefix is the exact prefix of the same \(T\) by (1)--(3).  The positivity
and goodness premises for Lemma 9n are fields of the
\noderef{HierarchyCounterexample} \(P.\mathit{counterexample}\), which is
one of the certified objects supplied by
\noderef{HierarchyReplacementPresentation}.

Put
\[
                        i=j-2,\qquad m=|A|.
\]
Since \(2\le j\), natural-number subtraction gives \(j=i+2\).  Since
\(j<\ell\), we also have \(i<\ell\).  Apply \noderef{NineN} to \(P\)'s
replacement ETT, replacement base, hierarchy \(H\), good-hierarchy
witness, the elementarity just obtained, and the
\(\mathit{largeMissingSet}\) field of
\(P.\mathit{minimumContext}\), at the index \(i\).  By (3), the set
called \(T(y_i)\) in that lemma is exactly \(B\).  In the notation of the
present statement, its old tail-color set is \(U_i\), and its available
set is
\[
 X=(\overline\mu(B)\setminus M_0)\setminus U_i.              \tag{4}
\]
Consequently Lemma 9n gives
\[
 |X|\ge
 2\,P.\mathit{corresponding.replacementETT.rounds}+11,        \tag{5}
\]
and, if
\[
                      |X\setminus(\Gamma\cup D)|\le4,         \tag{6}
\]
it gives an injectively indexed family
\((\eta_h)_{h\in\operatorname{Fin}7}\) such that, with
\(Q_h=\Gamma^q_{\eta_h}\cup\{\eta_h\}\),
\[
 \eta_h\in D\cap\overline\mu(B)\quad\hbox{and}\quad
                         Q_h\cap U_i=\varnothing              \tag{7}
\]
for every \(h\).

Only two new exact edge positions occur between \(U_i\) and \(U_j\).  To
record this without identifying parallel edges, set
\[
                              R=U_j\setminus U_i.
\]
Take \(c\in R\), and choose the index \(r\) in the definition of \(U_j\).
If \(r<m+i\), the very same index and the very same exact edge
\(T.\mathit{edge}(r)\) put \(c\) in \(U_i\), a contradiction.  Hence
\(m+i\le r<m+j=m+i+2\), so
\[
                         r=m+i\quad\hbox{or}\quad r=m+i+1.     \tag{8}
\]
The inequality \(j<\ell\), (2), and the fact that \(A\) is an ordered
initial segment of \(T\) from \noderef{TashkinovHierarchy} ensure that
both numbers in (8) are legitimate positions whenever they occur.
Thus \(R\) is contained in the set of colors assigned by \(\mu\) to these
two individual edge identities, and therefore
\[
                              |R|\le2,\qquad U_j\subseteq U_i\cup R. \tag{9}
\]

We split according to (6).  First suppose
\(|X\setminus(\Gamma\cup D)|>4\).  The set
\(R\cup\{\alpha\}\) has cardinality at most three by (9), whereas
\(X\setminus(\Gamma\cup D)\) has at least five members.  Choose
\[
 \beta\in (X\setminus(\Gamma\cup D))\setminus(R\cup\{\alpha\}). \tag{10}
\]
Then \(\beta\ne\alpha\), and (4) gives
\(\beta\in\overline\mu(B)\), \(\beta\notin M_0\), and
\(\beta\notin U_i\).  Equation (10) also gives
\(\beta\notin\Gamma\), \(\beta\notin D\), and \(\beta\notin R\).
The containment in (9) now yields \(\beta\notin U_j\), while the two
nonmembership conclusions give \(\beta\notin M_0\cup D\).  This is the
first required alternative.

It remains to assume (6), and take the seven colors and blocks in (7).
We first check explicitly that these seven blocks are pairwise disjoint.
By \noderef{HierarchyStoredColorsAt}, a stored color in \(D\) is not
missing on \(A\).  On the other hand, the \(\mathit{avoidance}\) field of
\noderef{GoodHierarchy}, interpreted using
\noderef{HierarchyStageMissingColors}, puts every member of every
\(\Gamma^q_\eta\), for \(\eta\in D\), among the colors missing on the
last relevant hierarchy stage.  If \(q>0\), that stage is \(A\) by
\noderef{HierarchyStarredStage}; if \(q=0\), it is \(T_n\), which is an
ordered initial segment of \(A\) by \noderef{HierarchyBase} and
\noderef{TashkinovHierarchy}.  In either case
\noderef{TreeMissingColors} shows that every Gamma color is missing on
\(A\).  Hence no member of \(D\) belongs to any Gamma pair.             \tag{11}

For distinct \(h,h'\), injectivity gives
\(\eta_h\ne\eta_{h'}\).  The \(\mathit{disjointPairs}\) field of
\noderef{GoodHierarchy} makes their two Gamma pairs disjoint; (11) removes
both possible stored-color/Gamma cross-intersections, and the two stored
colors themselves are distinct.  Therefore
\[
                              Q_h\cap Q_{h'}=\varnothing.       \tag{12}
\]

Let \(I_R\) be the indices \(h\) for which \(Q_h\cap R\ne\varnothing\).
For each \(h\in I_R\), choose one color in this intersection.  By (12),
the chosen colors for two distinct indices are distinct.  This gives an
injection \(I_R\hookrightarrow R\), so (9) implies \(|I_R|\le2\).
Likewise, (12) shows that at most one index \(h\) has
\(\alpha\in Q_h\).  Thus at most three of the seven indices are forbidden
by meeting \(R\) or containing \(\alpha\).  Choose an index \(h_0\) which
is not forbidden, and set
\[
                              \beta=\eta_{h_0}.                 \tag{13}
\]
By (7), \(\beta\in D\cap\overline\mu(B)\) and
\(Q_{h_0}\cap U_i=\varnothing\).  The choice of \(h_0\) gives
\(\alpha\notin Q_{h_0}\), hence \(\beta\ne\alpha\), and also
\(Q_{h_0}\cap R=\varnothing\).  Combining this with (9) and (7) gives
\[
       (\Gamma^q_\beta\cup\{\beta\})\cap U_j=\varnothing.      \tag{14}
\]
In particular \(\beta\notin U_j\).  Finally, (11) and \(\beta\in D\)
give \(\beta\notin\Gamma\).  Equations (13)--(14) are exactly the stored
second alternative and all the common conclusions.  The two cases prove
the lemma.
\end{proof}
